% ECONOMICS_PROBLEM: {"id": "H24-377", "title": "Inheritance taxation and dynastic wealth", "jel_code": "H24", "source_id": "OP-0050"}
% !TeX program = xelatex
\documentclass[11pt,a4paper]{article}
\usepackage[margin=23mm]{geometry}
\usepackage{fontspec}
\setmainfont{Latin Modern Roman}
\usepackage{amsmath,amssymb,mathtools}
\usepackage{xcolor,enumitem,booktabs,longtable,array,xurl,hyperref}
\definecolor{muted}{HTML}{53636C}
\hypersetup{colorlinks=true,urlcolor=blue,linkcolor=blue}
\setlength{\parindent}{0pt}
\setlength{\parskip}{5pt}
\newcommand{\R}{\mathbb R}
\newcommand{\E}{\mathbb E}
\newcommand{\N}{\mathbb N}
\newcommand{\ind}{\mathbf 1}
\newcommand{\Var}{\operatorname{Var}}
\newcommand{\Cov}{\operatorname{Cov}}
\newcommand{\supp}{\operatorname{supp}}
\DeclareMathOperator*{\argmin}{arg\,min}
\DeclareMathOperator*{\argmax}{arg\,max}
\newcommand{\entrymeta}[3]{{\sffamily\small\color{muted}\textbf{\texttt{#1}}\quad|\quad #2\par #3\par}}
\newcommand{\Problemref}[1]{\texttt{#1}}
\newcommand{\JELref}[2]{\texttt{#2} (JEL \texttt{#1})}
\begin{document}
\section*{Inheritance taxation and dynastic wealth}
\textbf{Problem ID:} \texttt{H24-377}\quad \textbf{JEL:} \texttt{H24}\par
% BEGIN ECONOMICS STATEMENT
\label{problem:OP-0050}
\label{atlas:prob:P10}
\noindent\textbf{Original atlas identifier:} P10; entry 50. \textbf{Field:} Public finance and taxation. \textbf{Original tractability label:} Hard.

\noindent\textbf{Classification:} Parameterized theoretical/empirical research agenda; not a certified unsolved theorem.

\subsection*{Definitions and assumptions}
Consider generations $g=0,\ldots,G$ linked by genealogical records and an initial joint distribution of assets, productivity, longevity, and bequest preferences. A policy $p$ specifies estate, recipient-inheritance, and inter vivos gift schedules, exemptions, valuation, deferrals, and enforcement; these are distinct instruments. Each dynamic model $m$ specifies saving, family-business investment, transfer timing, migration, mortality risk, and intentional or accidental bequests, together with equilibrium selection. Let $EV_{ig}(p;m)$ be monetary equivalent variation for beneficiary $i$ in generation $g$, including taxes and compliance costs, expressed in a common monetary numeraire before the stated generational welfare discount. For each generation $g$, let $\mathcal I_g$ be a fixed finite set of potential beneficiaries; set $EV_{ig}=0$ when a modeled birth/survival state makes $i$ absent, and declare the joint model law of these states and economic shocks. Fix discount factor $\beta\in(0,1]$, nonnegative generation/person weights $w_{ig}$, and present-value net revenue $N(p;m)$. Assume finite welfare sums and specify treatment of generations beyond $G$.

\subsection*{Formal research question}
\emph{Editorial formulation: the mathematical target below is not a theorem quoted from the references.}

For a legal feasible set $\mathcal P$ and a required present-value revenue target $\bar R$, characterize
\[
 \sup_{p\in\mathcal P:\ N(p;m)\geq\bar R}
 \sum_{g=0}^{G}\beta^g\sum_{i\in\mathcal I_g}w_{ig}\mathbb E_m[EV_{ig}(p;m)].
\]
Determine what linked multigenerational data and mortality/tax variation distinguish accidental bequests, warm-glow motives, altruism, and strategic exchange with heirs. For observationally equivalent motive models, identify the range of welfare gains from jointly taxing gifts and inheritances, allowing tax-base shifting and business valuation. A transfer's donor and recipient welfare must be evaluated jointly; avoid treating the same transfer as an aggregate resource gain twice. Report $\varepsilon$-optimal implementable rules or robust rankings when motive parameters are only set identified.

\subsection*{Interpretation and scope}
Deaths do not automatically provide exogenous bequest variation: health shocks can alter saving, caregiving, and earnings before death. Family firms require a liquidity and investment model because valuation discounts and deferral may change both genuine production and legal relabeling. A dynasty's objective is not the planner's objective unless an explicit ethical assumption equates them. Generation weights determine how much redistribution across heirs or cohorts is valued; those weights cannot be recovered from estate records. Mobility changes fiscal exposure and possibly the relevant social population, so keep the latter explicit. The research target is joint motive identification and feasible tax design over long horizons, with existing sufficient-statistic inheritance-tax benchmarks treated as established results.

\subsection*{Source audit and status}
[S1] is a review and [S2] derives inheritance-tax formulas. The bare Saez--Zucman (2019) citation is underdetermined; [S3] is explicitly selected as a contextual book, not asserted to be the unique intended work. These sources do not certify the proposed multigenerational extension's 2026 status.

\subsection*{References}
\begin{enumerate}[label={[S\arabic*]},leftmargin=*,itemsep=0.6em]
\item Wojciech Kopczuk (2013). Taxation of Intergenerational Transfers and Wealth. In Alan J. Auerbach, Raj Chetty, Martin Feldstein, and Emmanuel Saez (eds.), Handbook of Public Economics, vol. 5, chapter 6, 329--390. Elsevier. DOI: 10.1016/B978-0-444-53759-1.00006-6.
\url{https://www.sciencedirect.com/science/article/pii/B9780444537591000066}
\newline\emph{Locator and support limits:} Chapter review: transfer taxation and heterogeneous motives; review is not proof that the proposed extension remains unresolved.
\item Thomas Piketty and Emmanuel Saez (2013). A Theory of Optimal Inheritance Taxation. Econometrica 81(5), 1851--1886. DOI: 10.3982/ECTA10712.
\url{https://onlinelibrary.wiley.com/doi/abs/10.3982/ECTA10712}
\newline\emph{Locator and support limits:} Abstract and model: inheritance-tax sufficient statistics; identifiable dynastic dynamics require additional data and restrictions.
\item Emmanuel Saez and Gabriel Zucman (2019). The Triumph of Injustice: How the Rich Dodge Taxes and How to Make Them Pay. W. W. Norton. ISBN 9781324002727.
\url{https://wwnorton.co.uk/books/9781324002727-the-triumph-of-injustice}
\newline\emph{Locator and support limits:} Publisher title and authors: broad tax-policy synthesis. UK publication date is 12 November 2019.
\end{enumerate}

\subsection*{Common conventions: Source mathematical conventions}

Unless an entry states otherwise, agent and firm sets are finite. State and action spaces are measurable, strategies and outcome maps are measurable, probability laws are specified, and expectations exist for the targets under discussion. Infinite-horizon discount factors lie in $(0,1)$ and discounted objectives are finite. Norms, tolerances and complexity input sizes must be specified for computational comparisons. Constants or primitives that appear locally are inputs, not empirical facts asserted by this catalogue.

\textbf{Identification.} A statistical model $m\in\mathcal M$ implies an observable law $P_m$ and target $\theta(m)$. At law $P$, its identified set is
\[
\Theta(P;\mathcal M)=\{\theta(m):m\in\mathcal M,\ P_m=P\}.
\]
Point identification holds when this set is a singleton, or equivalently when $P_{m_1}=P_{m_2}$ implies $\theta(m_1)=\theta(m_2)$ for all compatible models. Sharp bounds mean bounds attained, or approached when the boundary is not attained, by compatible models. Writing this set is a definition, not a solution: the substantive task is to establish informative restrictions and derive or estimate the set. An unrestricted-confounding model may give only trivial bounds.

\textbf{Inference.} Identification, estimability and valid uncertainty quantification are separate questions. Fix $\alpha\in(0,1)$. A confidence procedure $C_n$ over a declared law class $\mathcal P$ has asymptotic uniform coverage at level $1-\alpha$ when
\[
\liminf_{n\to\infty}\inf_{P\in\mathcal P}\Pr_P\{\theta(P)\in C_n\}\geq1-\alpha.
\]
For set-identified targets the coverage event must be defined separately, for example coverage of the full identified set. If an entry uses identification-indexed models instead of uniquely indexed laws, the same coverage convention is applied over those models. Pointwise limits do not establish uniform validity, and asymptotic validity does not guarantee finite-sample precision. Sample sizes, dependence structures and weak-identification sequences are part of each application.

\textbf{Counterfactuals.} $Y_i(a)$ is a potential outcome under a specified intervention, including its timing, implementation and versions. It is not $Y_i$ conditional on voluntarily choosing $a$. A full intervention vector permits interference; shorthand scalar treatment notation does not silently impose no interference. Assignment, selection, attrition, anticipation and measurement must be specified. A claim about an instrument's local effect is not automatically a population policy effect.

\textbf{Economic equilibrium.} A counterfactual model includes best responses, constraints, feasibility, market clearing, state transitions and expectations consistent with the stated information. When several equilibria exist, either a declared selector is an input or the target is set-valued. Comparisons across policy regimes must use the stated selection convention. Reduced-form relationships are not presumed invariant after an equilibrium-changing intervention.

\textbf{Optimization and welfare.} Optimization is over the stated feasible set. If attainment is not established, $\sup$ or $\inf$ and $\varepsilon$-optimal choices replace an asserted maximizer or minimizer. For example, with value $V=\sup_{a\in\mathcal A}W(a)<\infty$, an $\varepsilon$-optimal set is $\{a\in\mathcal A:W(a)\geq V-\varepsilon\}$ for $\varepsilon>0$. Benefits, costs, risk and health or environmental outcomes can be aggregated only using declared units and conversion weights. Interpersonal utility weights, the treatment of fiscal transfers and foreign profits, discounting, and the behavioral welfare criterion are explicit inputs. A comparative-static derivative additionally requires existence and appropriate differentiability of the selected counterfactual.

\textbf{Sources.} Local labels [S1], [S2], and so forth restart in every question. Locators identify the relevant abstract, model, section or result at the level actually checked; a bibliographic or abstract-level verification is not represented as inspection of an entire proof. Working papers are distinguished from later publications and changing versions. Source summaries are paraphrased. All URL checks and editorial formulations refer to the audit date stated above.
% END ECONOMICS STATEMENT
\end{document}
